\section{Exercise 3 : Application of the principle of maximum entropy}

\begin{frame}{Exercise 3: Application of the principle of maximum entropy}

This exercise introduces the \textbf{principle of maximum entropy}, a method used to determine the probability distribution of random variables based on limited information—specifically, subject to a known constraint, which in this case is their \textbf{average value}. When only the mean value of a system is known, the principle of maximum entropy suggests that we should choose the probability distribution that maximizes uncertainty (or entropy), reflecting the most unbiased representation of the data.

To solve this problem, one effective approach is the use of \textbf{Lagrange optimization}, a mathematical technique that allows us to maximize the entropy while respecting the constraint of the mean.

The exercise consists of two parts:

\begin{itemize}
    \item The first part involves a random variable with \textbf{two outcomes}
    \item The second part expands the problem by introducing additional outcomes, increasing the number of possible values for the random \textbf{variable to four}.
\end{itemize}

\end{frame}

\begin{frame}{Data Selection}

For this exercise, we chose a real-world dataset with integer outcomes: the occurrence times of earthquakes in Japan from recent dates. You can find more information about on this \href{https://earthquaketrack.com/p/japan/recent}{link}.

\begin{itemize}
    \item \textbf{For the First Part our data are } : 3 earthquakes with magnitude above 4, and 6 with magnitude lower than 4 on 14th of October 2024 in Japan.
    \item \textbf{For the Second Part our data are } : 1 earthquakes with magnitude above 5, and 5 times with magnitude lower than 3 on 12th of October 2024 in Japan. 
\end{itemize}  
\end{frame}
\begin{frame}{Lagrange Optimization}

Lets consider our constraints that we are working on to find the probability distribution: 

\begin{itemize}
    \item The normalization constraint : \(\sum_{i} p_i  = 1 \)
    \item The mean values : \( \mu = \sum_{i} p_i x_i \)
\end{itemize}

By applying the entropy equation under above constraint in Lagrange optimization, we can find the probability distribution by:

\[ P(x_i) = \frac{\beta^{x_i}}{\sum_{i=1}^{n} \beta^{x_i}} \]

Where the \( \beta \) is defined based on the this formula: 

\[ \sum_{i=1}^{n} \beta^{x_i} x_i = \mu (\sum_{i=1}^{n} \beta^{x_i})  \]

\end{frame}

\begin{frame}{Two Outcomes Example}

As we defined our samples data we can write our alphabet \( \Omega_x \) based on recorded of earthquake frequencies that is \( \Omega_x = [3,6] \). Therefore, it forms this equation:

\[ x_0 \beta^{x_0} + x_1 \beta^{x_1} = \mu( \beta^{x_0} + \beta^{x_1} ) \]  

For performing our experiments of finding the probabilities distribution, we calculated with 3 different mean values ( \(\mu\) ) as each of them are selected based on arithmetic average of random variables. where the average is equal to 4.5.

\begin{itemize}
    \item Lets assume \(\mu= 5 \) \textbf{(above the arithmetic average)}
    \item Therefore the \(\beta \approx 1.259 \).
    \item Then the probabilities for each outcomes are: \( P(x_0 = 3) = 0.333, P(x_1 = 6) = 0.667\) 
\end{itemize}
    
\end{frame}
\begin{frame}{Two outcomes example with different mean values}

If we set the \(\mu = 4.5\) as \textbf{equal} to the arithmetic average, therefore we have:
\begin{itemize}
    \item the \(\beta = 1\) and the probabilities of two outcomes becomes: \( P(x_0 = 1) = P(x_1 = 3) = 0.5 \)
\end{itemize}

And if ser the \(\mu = 4\)  \textbf{less than} the arithmetic average, therefore we will have:

\begin{itemize}
    \item the \(\beta \approx 0.793\) and the probabilities of two outcomes becomes: \( P(x_0 = 3) = 0.667, P(x_1 = 6) = 0.333 \)
\end{itemize}
    
\end{frame}

\begin{frame}{Four outcomes Example}

To extend the impacts of using different mean values, this exercise moves forward to experiment with the same process on four outcomes.

Now The alphabet is \(\Omega_x = [3,6,1,5]\). 

For this scenario, we calculated the probabilities based on different mean values that are listed as below.

\begin{itemize}
    \item \textbf{ \(\mu = 5\) Higher than Average}(with \(\beta \approx 1.5436\)) : [
    \( P(x_0 = 3) = 0.1337, P(x_1 = 6) = 0.4917, P(x_2 = 1) = 0.0561, P(x_3 = 5) = 0.3185 \)]
    \item \textbf{ \(\mu = 3.75\) Equal to Average:} [
    \( P(x_0 = 3) = P(x_1 = 6) = P(x_2 = 1) = P(x_3 = 5) = 0.5 \)]
    \item \textbf{ \(\mu = 2\) Less than Average}(with \(\beta \approx 0.5903\)) : [ 
    \( P(x_0 = 3) = 0.2261, P(x_1 = 6) = 0.0465 , P(x_2 = 1) = 0.6486, P(x_3 = 5) = 0.0788 \)]
\end{itemize}
    
\end{frame}
\begin{frame}{Results}

We are representing the difference between fixing different mean values by plotting them alongside each other to make better comparison.

\begin{figure}
    \includegraphics[width=1.0\linewidth]{images/Ex3all3.png}
    \label{fig:enter-label}
\end{figure}
    
\end{frame}

\begin{frame}{Discussion}
\begin{itemize}
    \item \textbf{Mean equal to the average:} The probability distribution was balanced, maximizing uncertainty with both outcomes having similar probabilities.
    \item \textbf{Mean above the average:} The distribution skewed toward the higher outcome, reducing entropy as the higher value became more likely.
    \item \textbf{Mean below the average:} The distribution shifted toward the lower outcome, similarly reducing entropy by favoring the smaller value.
\end{itemize}
	
In General, As the mean moves away from the arithmetic average, the probability distribution becomes more skewed, reducing entropy. The system adjusts to reflect the bias introduced by the higher or lower mean, but it still adheres to the principle of maximum entropy by maximizing uncertainty given the new constraints. This illustrates how the maximum entropy method adapts to different known values while maintaining as much uncertainty as possible.
    
\end{frame}